% !TEX root = ../main.tex
\chapter{Maxwell-to-Signal Derivation (Field Side)}
\label{app:2dfts_maxwell}

This appendix gives the \emph{field-side} derivation for coherent 2D-FTS/2DES. Starting from Maxwell's equations, it shows how the radiated, heterodyne-detected signal field is linear in the third-order polarisation $\boldsymbol{P}^{(3)}$. This provides the proportionality used in~\autoref{sec:spctr_pol_and_macr_response}.

The argument proceeds in three steps. First, Maxwell's equations lead to a driven wave equation in which the macroscopic polarisation acts as the source. Second, under the slowly varying envelope approximation, the emitted signal field follows the corresponding component of $\boldsymbol{P}^{(3)}$. 
The microscopic expression for $\boldsymbol{P}^{(3)}$ is derived on the matter side in~\autoref{app:response_functions}.

\section{From Maxwell's Equations to a Wave Equation for the Radiation Field}
\label{app:sec:maxwell_to_wave}

This section makes explicit why the macroscopic polarisation is the source of the emitted optical field~\cite{mancal2022nonlinearopticalspectroscopy}. The central point is that only the \emph{transverse} part of the electromagnetic field propagates as radiation, and it is sourced by the transverse current. In condensed-matter optics, that current is naturally written in terms of the polarisation.

\subsection{Maxwell's Equations and Electromagnetic Potentials}
\label{app:subsec:maxwell_eqs}
In natural units ($\hbar = c = 1$) and Heaviside--Lorentz form, Maxwell's equations in the presence of charge and current densities $\varrho(\boldsymbol{r},t)$ and $\boldsymbol{j}(\boldsymbol{r},t)$ read \cite{guidry1991gaugefieldtheories}
\begin{align}
	\boldsymbol{\nabla} \times \boldsymbol{B}(\boldsymbol{r},t)
	 & = \boldsymbol{j}(\boldsymbol{r},t) + \frac{\partial \boldsymbol{E}(\boldsymbol{r},t)}{\partial t},
	\label{eq:app_maxwell_ampere}
	\\
	\boldsymbol{\nabla} \times \boldsymbol{E}(\boldsymbol{r},t)
	 & = -\frac{\partial \boldsymbol{B}(\boldsymbol{r},t)}{\partial t},
	\label{eq:app_maxwell_faraday}
	\\
	\boldsymbol{\nabla}\cdot \boldsymbol{E}(\boldsymbol{r},t)
	 & = \varrho(\boldsymbol{r},t),
	\label{eq:app_maxwell_gaussE}
	\\
	\boldsymbol{\nabla}\cdot \boldsymbol{B}(\boldsymbol{r},t)
	 & = 0.
	\label{eq:app_maxwell_gaussB}
\end{align}
Here $\boldsymbol{E}$ is the electric field and $\boldsymbol{B}$ the magnetic field.

It is convenient to express $\boldsymbol{E}$ and $\boldsymbol{B}$ in terms of a vector potential $\boldsymbol{A}(\boldsymbol{r},t)$ and a scalar potential $\phi(\boldsymbol{r},t)$,
\begin{equation}
	\boldsymbol{B} = \boldsymbol{\nabla}\times \boldsymbol{A},
	\qquad
	\boldsymbol{E} = -\frac{\partial \boldsymbol{A}}{\partial t} - \boldsymbol{\nabla}\phi.
	\label{eq:app_potentials_def}
\end{equation}
This automatically satisfies Faraday's law \cref{eq:app_maxwell_faraday} and Gauss' law for magnetism \cref{eq:app_maxwell_gaussB}, reducing the number of independent equations.

\subsection{Gauge Freedom and Coulomb Gauge}
\label{app:subsec:gauge_freedom}
To proceed, a gauge must be chosen. The potentials in~\autoref{eq:app_potentials_def} are not unique: for any scalar function $\chi(\boldsymbol{r},t)$,
\begin{equation*}
	\boldsymbol{A}' = \boldsymbol{A} + \boldsymbol{\nabla}\chi,
	\qquad
	\phi' = \phi - \frac{\partial \chi}{\partial t},
\end{equation*}
leaves $\boldsymbol{E}$ and $\boldsymbol{B}$ invariant. For optics, the Coulomb gauge is used,
\begin{equation*}
\boldsymbol{\nabla}\cdot\boldsymbol{A} = 0.
\end{equation*}
Inserting \autoref{eq:app_potentials_def} into Gauss' law \autoref{eq:app_maxwell_gaussE} then gives the Poisson equation
\begin{equation*}
	\boldsymbol{\nabla}^2\phi(\boldsymbol{r},t) = -\varrho(\boldsymbol{r},t),
\end{equation*}
so $\phi$ is determined instantaneously in space. This does not violate causality; it reflects the gauge-dependent description of the Coulomb interaction. The propagating optical field resides in the transverse sector, discussed next.

\subsection{Transverse/Longitudinal Decomposition}
\label{app:subsec:transverse_longitudinal}
To separate propagating radiation from the Coulomb part, decompose fields into transverse and longitudinal components. Only the transverse component carries radiative degrees of freedom.

Any sufficiently well-behaved vector field $\boldsymbol{F}(\boldsymbol{r})$ admits the Helmholtz decomposition,
\begin{equation*}
	\boldsymbol{F} = \boldsymbol{F}^{\perp} + \boldsymbol{F}^{\parallel},
\end{equation*}
where the transverse part is divergence-free and the longitudinal part is curl-free,
\begin{equation*}
	\boldsymbol{\nabla}\cdot\boldsymbol{F}^{\perp}=0,\qquad \boldsymbol{\nabla}\times\boldsymbol{F}^{\parallel}=0.
\end{equation*}
Applying this decomposition to the electric field in~\autoref{eq:app_potentials_def},
\begin{equation*}
	\boldsymbol{E} = \boldsymbol{E}^{\perp} + \boldsymbol{E}^{\parallel},
\end{equation*}
with $\boldsymbol{E}^{\perp} = -\partial_t \boldsymbol{A}$ in Coulomb gauge and $\boldsymbol{E}^{\parallel} = -\boldsymbol{\nabla}\phi$. Maxwell's equations \cref{eq:app_maxwell_ampere,eq:app_maxwell_faraday,eq:app_maxwell_gaussE,eq:app_maxwell_gaussB} then split into independent transverse and longitudinal parts,
\begin{align*}
	\boldsymbol{\nabla}\times\boldsymbol{E}^{\perp}     & = -\frac{\partial \boldsymbol{B}}{\partial t},
	                                                    &
	\boldsymbol{\nabla}\times\boldsymbol{E}^{\parallel} & = 0,
	\\
	\boldsymbol{\nabla}\cdot\boldsymbol{E}^{\parallel}  & = \varrho,
	                                                    &
	\boldsymbol{\nabla}\cdot\boldsymbol{E}^{\perp}      & = 0.
\end{align*}
Faraday's law constrains only the transverse field because $\boldsymbol{\nabla}\times\boldsymbol{E}^{\parallel}=0$ identically, while Gauss' law fixes the longitudinal field through the charge density. The propagating and Coulomb degrees of freedom are therefore separated.

\medskip

To understand which current sources what, recall that charge and current are not independent: they are linked by the continuity equation
\begin{equation}
	\frac{\partial \varrho(\boldsymbol{r},t)}{\partial t} + \boldsymbol{\nabla}\cdot\boldsymbol{j}(\boldsymbol{r},t)=0.
	\label{eq:app_continuity}
\end{equation}
Decomposing $\boldsymbol{j}=\boldsymbol{j}^{\perp}+\boldsymbol{j}^{\parallel}$ and using $\boldsymbol{\nabla}\cdot\boldsymbol{j}^{\perp}=0$ shows that \autoref{eq:app_continuity} constrains only the longitudinal part,
\begin{equation*}
	\frac{\partial \varrho}{\partial t} + \boldsymbol{\nabla}\cdot\boldsymbol{j}^{\parallel}=0.
\end{equation*}
In this sense, $\boldsymbol{j}^{\parallel}$ is tied to Coulomb near-field physics, while radiative fields are sourced by $\boldsymbol{j}^{\perp}$.

Two immediate consequences follow.
\begin{itemize}
	\item Because $\boldsymbol{\nabla}\cdot\boldsymbol{B}=0$, the magnetic field is purely transverse, so $\boldsymbol{B}=\boldsymbol{B}^{\perp}$.
	\item In Coulomb gauge $\boldsymbol{\nabla}\cdot\boldsymbol{A}=0$, the vector potential is purely transverse, so $\boldsymbol{A}=\boldsymbol{A}^{\perp}$.
\end{itemize}
Combining the decomposition with Coulomb gauge and \autoref{eq:app_potentials_def} yields
\begin{equation*}
	\boldsymbol{E}^{\perp} = -\frac{\partial \boldsymbol{A}}{\partial t},
	\qquad
	\boldsymbol{E}^{\parallel} = -\boldsymbol{\nabla}\phi.
\end{equation*}
The radiation field relevant for propagating optics is $\boldsymbol{E}^{\perp}$. The longitudinal field $\boldsymbol{E}^{\parallel}$ is fixed instantaneously by $\varrho$ via the Poisson equation and does not describe propagation.

\subsection{Wave Equation for the Transverse Electric Field}
\label{app:subsec:wave_equation}
With the radiative field isolated, the wave equation it satisfies is derived. Starting from Faraday's law \autoref{eq:app_maxwell_faraday}, take the curl of both sides,
\begin{equation*}
	\boldsymbol{\nabla}\times(\boldsymbol{\nabla}\times\boldsymbol{E})
	= -\frac{\partial}{\partial t}\big(\boldsymbol{\nabla}\times\boldsymbol{B}\big).
\end{equation*}
Next, eliminate $\boldsymbol{\nabla}\times\boldsymbol{B}$ using Ampère--Maxwell's law \autoref{eq:app_maxwell_ampere} to obtain
\begin{equation*}
	-\,\boldsymbol{\nabla}\times(\boldsymbol{\nabla}\times\boldsymbol{E})
	= \frac{\partial^2\boldsymbol{E}}{\partial t^2} + \frac{\partial \boldsymbol{j}}{\partial t}.
\end{equation*}
Only transverse parts contribute to the curl--curl operator because $\boldsymbol{\nabla}\times\boldsymbol{E}^{\parallel}=0$. Likewise, the radiated field is sourced by the transverse current $\boldsymbol{j}^{\perp}$. Projecting onto the transverse sector gives
\begin{equation*}
	-\,\boldsymbol{\nabla}\times(\boldsymbol{\nabla}\times\boldsymbol{E}^{\perp})
	= \frac{\partial^2\boldsymbol{E}^{\perp}}{\partial t^2} + \frac{\partial \boldsymbol{j}^{\perp}}{\partial t}.
\end{equation*}
Using $\boldsymbol{\nabla}\times(\boldsymbol{\nabla}\times\boldsymbol{E}^{\perp}) = \boldsymbol{\nabla}(\boldsymbol{\nabla}\cdot\boldsymbol{E}^{\perp}) - \boldsymbol{\nabla}^2\boldsymbol{E}^{\perp}$ and $\boldsymbol{\nabla}\cdot\boldsymbol{E}^{\perp}=0$ yields the driven wave equation
\begin{equation}
	\boldsymbol{\nabla}^2\boldsymbol{E}^{\perp}(\boldsymbol{r},t)
	- \frac{\partial^2\boldsymbol{E}^{\perp}(\boldsymbol{r},t)}{\partial t^2}
	= \frac{\partial \boldsymbol{j}^{\perp}(\boldsymbol{r},t)}{\partial t}.
	\label{eq:app_wave_E_perp_j}
\end{equation}

\subsection{From Transverse Current to Polarisation}
\label{app:subsec:current_to_polarisation}
The wave equation \autoref{eq:app_wave_E_perp_j} is expressed in terms of the current $\boldsymbol{j}^{\perp}$. In optics it is more convenient to work with the macroscopic polarisation $\boldsymbol{P}(\boldsymbol{r},t)$, which describes the motion of bound charges in a neutral dielectric on coarse-grained scales. The associated current equals the time derivative of the polarisation,
\begin{equation}
	\boldsymbol{j}(\boldsymbol{r},t) = \frac{\partial \boldsymbol{P}(\boldsymbol{r},t)}{\partial t}.
	\label{eq:app_j_equals_Pdot}
\end{equation}
Together with the continuity equation \autoref{eq:app_continuity}, this implies the standard relation between bound charge density and polarisation,
\begin{equation*}
	\varrho_b(\boldsymbol{r},t) = -\boldsymbol{\nabla}\cdot\boldsymbol{P}(\boldsymbol{r},t).
\end{equation*}
Substituting the polarisation current \autoref{eq:app_j_equals_Pdot} into \autoref{eq:app_wave_E_perp_j} gives
\begin{equation}
	\boldsymbol{\nabla}^2\boldsymbol{E}^{\perp}(\boldsymbol{r},t)
	- \frac{\partial^2\boldsymbol{E}^{\perp}(\boldsymbol{r},t)}{\partial t^2}
	= \frac{\partial^2\boldsymbol{P}(\boldsymbol{r},t)}{\partial t^2}.
	\label{eq:app_wave_with_P}
\end{equation}
\autoref{eq:app_wave_with_P} shows directly that the macroscopic polarisation is the source of the radiated transverse field. Measuring the emitted optical field therefore probes the polarisation induced by the incident light.

\medskip

For nonlinear spectroscopy, write $\boldsymbol{P}=\boldsymbol{P}^{(1)}+\boldsymbol{P}_{\mathrm{NL}}$, where $\boldsymbol{P}_{\mathrm{NL}}$ collects higher-order terms such as $\boldsymbol{P}^{(3)}$. The linear response $\boldsymbol{P}^{(1)}$ is absorbed into an effective refractive index $n(\omega)$ with $n^2 = 1 + \chi^{(1)}$. With this, \autoref{eq:app_wave_with_P} becomes
\begin{equation}
	\boldsymbol{\nabla}^2 \boldsymbol{E}(\boldsymbol{r},t)
	- n^2\,\frac{\partial^2}{\partial t^2}\boldsymbol{E}(\boldsymbol{r},t)
	= \frac{\partial^2}{\partial t^2}\boldsymbol{P}_{\mathrm{NL}}(\boldsymbol{r},t).
	\label{eq:app_wave_eq_nonlinear}
\end{equation}
In a three-pulse experiment operated in the weak-field regime, the dominant nonlinear source is the third-order term,
\begin{equation*}
	\boldsymbol{P}_{\mathrm{NL}}(\boldsymbol{r},t) \approx \boldsymbol{P}^{(3)}(\boldsymbol{r},t),
\end{equation*}
where $\boldsymbol{P}^{(3)}$ is cubic in the incident field.

The remaining task is to connect $\boldsymbol{P}^{(3)}$ to a measurable signal field using the SVEA.

% VERSION: (slab along the z direction)
\section{Slowly Varying Envelope Approximation and the Emitted Field}
\label{app:sec:svea}

The driven wave equation \autoref{eq:app_wave_eq_nonlinear} is solved to connect the nonlinear polarisation to the observed signal field. The slowly varying envelope approximation exploits the separation between the fast optical oscillation at carrier frequency $\omega_L$ and the slow evolution of envelopes. In this limit, the phase-matched signal field is proportional to the corresponding component of the nonlinear polarisation.

\subsection{Setup: Fourier Components and Ansatz}
\label{app:subsec:fourier_ansatz}

Consider a slab sample of thickness $L$ along $z$. The third-order polarisation contains Fourier components with different wavevector combinations $\boldsymbol{k}_\text{S}$. Phase matching selects which signal direction is observed, as discussed in~\autoref{subsec:spctr_phase_matching_condition}. A single component is written with a slowly varying envelope,
\begin{equation}
	\boldsymbol{P}_{\mathrm{NL}}(\boldsymbol{r},t) \supset \boldsymbol{P}_{\mathrm{S}}^{(3)}(t)\,e^{\mathrm{i}(\boldsymbol{k}_\text{S}\cdot\boldsymbol{r}-\omega_{L} t)} + \text{c.c.}
	\label{eq:app_pol_component}
\end{equation}
where $\boldsymbol{P}_{\mathrm{S}}^{(3)}(t)$ varies slowly in time. 
A corresponding forward-propagating field solution is sought with matching carrier phase,
\begin{equation}
	\boldsymbol{E}_{\mathrm{S}}(\boldsymbol{r},t) = \tfrac{1}{2}\,\boldsymbol{E}_{\mathrm{S}}(z,t)\,e^{\mathrm{i}(k'_{\mathrm{S}} z-\omega_{L} t)} + \text{c.c.},
	\label{eq:app_field_ansatz}
\end{equation}
with dispersion relation $k'_{\mathrm{S}}=n(\omega_{L})\omega_{L}$.

\subsection{SVEA Conditions}
\label{app:subsec:svea_conditions}

The SVEA assumes that the envelopes $\boldsymbol{E}_{\mathrm{S}}(z,t)$ and $\boldsymbol{P}_{\mathrm{S}}^{(3)}(t)$ vary slowly compared to the optical oscillation,
\begin{equation}
	\left|\partial_t \boldsymbol{E}_{\mathrm{S}}\right| \ll \omega_{L}\,\left|\boldsymbol{E}_{\mathrm{S}}\right|,
	\qquad
	\left|\partial_z^2 \boldsymbol{E}_{\mathrm{S}}\right| \ll k'_{\mathrm{S}}\,\left|\partial_z \boldsymbol{E}_{\mathrm{S}}\right|.
	\label{eq:app_svea_conditions}
\end{equation}
This scale separation allows higher-order derivatives of the envelope to be neglected when they appear alongside rapidly oscillating carrier factors.

\subsubsection{Spatial and Temporal Derivatives}
\label{app:subsubsec:spatial_temporal}
To isolate the envelope dynamics, keep only the positive-frequency part and suppress explicit time dependence momentarily. For a positive-frequency component
\begin{equation*}
	\boldsymbol{E}_{\mathrm{S}}^{(+)}(z) = \tfrac{1}{2}\,\boldsymbol{E}_{\mathrm{S}}(z)\,e^{\mathrm{i} k'_{\mathrm{S}} z},
\end{equation*}
one finds
\begin{align}
	\partial_z \boldsymbol{E}_{\mathrm{S}}^{(+)}
	 & = \tfrac{1}{2}\left(\partial_z \boldsymbol{E}_{\mathrm{S}} + \mathrm{i} k'_{\mathrm{S}}\boldsymbol{E}_{\mathrm{S}}\right)e^{\mathrm{i} k'_{\mathrm{S}} z},
	\notag\\
	\partial_z^2 \boldsymbol{E}_{\mathrm{S}}^{(+)}
	 & = \tfrac{1}{2}\left(\partial_z^2 \boldsymbol{E}_{\mathrm{S}} + 2 \mathrm{i} k'_{\mathrm{S}}\partial_z \boldsymbol{E}_{\mathrm{S}} - k_{\mathrm{S}}'^2\boldsymbol{E}_{\mathrm{S}}\right)e^{\mathrm{i} k'_{\mathrm{S}} z}
	\approx \tfrac{1}{2}\left(2 \mathrm{i} k'_{\mathrm{S}}\partial_z \boldsymbol{E}_{\mathrm{S}} - k_{\mathrm{S}}'^2\boldsymbol{E}_{\mathrm{S}}\right)e^{\mathrm{i} k'_{\mathrm{S}} z},
	\label{eq:app_second_derivative_svea}
\end{align}
where the last step uses \autoref{eq:app_svea_conditions} to neglect $\partial_z^2\boldsymbol{E}_{\mathrm{S}}$.

For the time dependence, the analogous approximation is
\begin{equation}
	\partial_t^2\!\left[\boldsymbol{E}_{\mathrm{S}}(z,t)e^{-\mathrm{i}\omega_{L} t}\right] \approx -\omega_{L}^2\,\boldsymbol{E}_{\mathrm{S}}(z,t)e^{-\mathrm{i}\omega_{L} t},
	\qquad
	\partial_t^2\!\left[\boldsymbol{P}_{\mathrm{S}}^{(3)}(t)e^{-\mathrm{i}\omega_{L} t}\right] \approx -\omega_{L}^2\,\boldsymbol{P}_{\mathrm{S}}^{(3)}(t)e^{-\mathrm{i}\omega_{L} t}.
	\label{eq:app_time_derivatives_svea}
\end{equation}
The carrier oscillation at $\omega_L$ dominates, so second derivatives of the envelopes can be dropped.

\subsection{Envelope Propagation Equation}
\label{app:subsec:envelope_propagation}

Insert \autoref{eq:app_pol_component} and \autoref{eq:app_field_ansatz} into \autoref{eq:app_wave_eq_nonlinear}. Using the SVEA derivative forms from \autoref{eq:app_second_derivative_svea} and \autoref{eq:app_time_derivatives_svea}, and imposing the dispersion relation $k_{\mathrm{S}}'^2=n^2(\omega_{L})\omega_{L}^2$, the leading carrier terms cancel. The result is a first-order propagation equation for the envelope,
\begin{equation}
	\mathrm{i} k'_{\mathrm{S}}\,\frac{\partial \boldsymbol{E}_{\mathrm{S}}(z,t)}{\partial z}
	= -\omega_{L}^2\,\boldsymbol{P}_{\mathrm{S}}^{(3)}(t)\,e^{\mathrm{i}\Delta k\,z},
	\label{eq:app_envelope_propagation}
\end{equation}
where $\Delta k = k_{\mathrm{S}} - k'_{\mathrm{S}}$ is the phase mismatch between the polarisation wavevector and the propagation constant $k'_{\mathrm{S}} = n(\omega_{L})\omega_{L}$. This equation shows how the field builds up through the sample, driven locally by the induced polarisation.

\subsection{Integration and Sinc Factor}
\label{app:subsec:integration_sinc}

\autoref{eq:app_envelope_propagation} is integrated across the sample. With no signal at the entrance, $\boldsymbol{E}_{\mathrm{S}}(0,t)=0$, integration from $z=0$ to $z=L$ gives the exiting signal field,
\begin{equation}
	\boldsymbol{E}_{\mathrm{S}}(L,t)
	= \mathrm{i}\,\frac{\omega_{L}}{n(\omega_{L})}\,\boldsymbol{P}_{\mathrm{S}}^{(3)}(t)\,L\,e^{\mathrm{i}\Delta k L/2}\,\mathrm{sinc}\!\left(\frac{\Delta k L}{2}\right).
	\label{eq:app_sinc}
\end{equation}
This result shows that the emitted signal field is directly proportional to the corresponding component of $\boldsymbol{P}_{\mathrm{S}}^{(3)}$, and the $\mathrm{sinc}$ term captures phase-mismatch suppression. The response is maximal near $\Delta k \approx 0$.

When the experiment is phase matched, $\Delta k\approx 0$, so $\mathrm{sinc}(\Delta k L/2)\to 1$ and $e^{\mathrm{i}\Delta k L/2}\to 1$. \autoref{eq:app_sinc} then reduces to
\begin{equation*}
	\boldsymbol{E}_{\mathrm{S}}(t)\equiv \boldsymbol{E}_{\mathrm{S}}(L,t)
	\propto \mathrm{i}\,\frac{\omega_{L}}{n(\omega_{L})}\,\boldsymbol{P}_{\mathrm{S}}^{(3)}(t)\,L.
\end{equation*}

In the frequency domain, this relation takes the analogue form
\begin{equation}
	\boldsymbol{E}_{\mathrm{S}}(\omega) = \kappa\, \mathrm{i}\omega\,\boldsymbol{P}^{(3)}(\omega),
	\label{eq:app_Es_iomegaP}
\end{equation}
where $\kappa$ is a real prefactor that depends on propagation conventions and experimental geometry. \autoref{eq:app_Es_iomegaP} is the Maxwell-side link used throughout: it makes the measurement a linear functional of $\boldsymbol{P}^{(3)}$.

\medskip

This completes the field-side derivation. The incident fields drive the matter system, producing $\boldsymbol{P}^{(3)}$ via the response-function formalism in~\autoref{app:response_functions}. That polarisation then acts as the source term that radiates the measured signal field. Detection models and the construction of 2D spectra from phase-selected time-domain signals are discussed in~\autoref{sec:spctr_2D_construction_formal}.